\documentclass[12pt]{article}
\usepackage{amsmath}
\usepackage{amssymb}
\usepackage{amsfonts}
\usepackage{amsthm}
\usepackage{bm}
\usepackage{array}
\usepackage[margin=1in]{geometry}
\usepackage[colorlinks=true,linkcolor=blue,citecolor=blue,urlcolor=blue]{hyperref}

\numberwithin{equation}{section}

\newtheorem{theorem}{Theorem}
\newtheorem{proposition}{Proposition}
\newtheorem{lemma}{Lemma}
\newtheorem{corollary}{Corollary}
\newtheorem{assumption}{Assumption}
\theoremstyle{definition}
\newtheorem{definition}{Definition}
\newtheorem{remark}{Remark}

\newcommand{\ind}{\mathbf 1}
\newcommand{\R}{\mathbb R}

\begin{document}

\begin{center}
{\Large\bf Math Logic File: Model, Lemmas, and Proofs}\\[6pt]
{\large Multiple Nuisance Breaks and Spurious Predictability under Persistent Regressors}
\end{center}

\section{Model Setup and Score Objects}

For a scalar array \(\bm a=(a_1,\ldots,a_T)'\), write
\[
    \|\bm a\|_T^2
    =
    T^{-1}\sum_{t=1}^T a_t^2,
    \qquad
    \|\bm a\|_\infty=\max_{t\leq T}|a_t|.
\]
Unless stated otherwise, \(O_p(\cdot)\) and \(o_p(\cdot)\) for growing
vectors and matrices are in Euclidean and spectral norm, respectively.

Let
\[
    \bm Y=(Y_1,\ldots,Y_T)',\qquad
    \bm X_{lag}=(X_0,\ldots,X_{T-1})',
\]
and let \(\bm w=(w_1,\ldots,w_T)'\), where \(w_t=w(X_{t-1})\) is a bounded saturated weight. The predictor satisfies
\begin{equation}
    X_t=\theta+\rho_TX_{t-1}+\varepsilon_t.
    \label{eq:predictor}
\end{equation}
The maintained outcome model is
\begin{equation}
    Y_t
    =
    \beta_0X_{t-1}
    +
    \sum_{j=0}^{q_0}
    m_j(W_{t-1})\ind\{k_{j,0}<t\leq k_{j+1,0}\}
    +
    U_t,
    \label{eq:model}
\end{equation}
where
\[
    0=k_{0,0}<k_{1,0}<\cdots<k_{q_0,0}<k_{q_0+1,0}=T.
\]
The coefficient \(\beta_0\) is stable over the full sample. The structural breaks are confined to the nonparametric nuisance component.

For a candidate partition \(\Lambda=(k_1,\ldots,k_q)\), write \(k_0=0\), \(k_{q+1}=T\), and
\[
    I_j(\Lambda)=\{t:k_j<t\leq k_{j+1}\}.
\]
The true partition is \(\Lambda_0=(k_{1,0},\ldots,k_{q_0,0})\). Define \(j_\Lambda(t)=j\) when \(t\in I_j(\Lambda)\). The true nuisance vector is
\[
    \mu_t=m_{j_{\Lambda_0}(t)}(W_{t-1}),
    \qquad
    \bm\mu=(\mu_1,\ldots,\mu_T)'.
\]

Let \(\bm P\) be the \(T\times K\) sieve matrix whose \(t\)-th row is
\[
    \bm P_K(W_{t-1})'=(p_1(W_{t-1}),\ldots,p_K(W_{t-1})).
\]
For a partition \(\Lambda\), let \(\bm D_j(\Lambda)\) be the diagonal matrix with diagonal entry \(\ind\{t\in I_j(\Lambda)\}\). Define
\begin{equation}
    \bm Q_{K,\Lambda}
    =
    \left[
    \bm D_0(\Lambda)\bm P,\,
    \bm D_1(\Lambda)\bm P,\,
    \ldots,\,
    \bm D_q(\Lambda)\bm P
    \right],
    \label{eq:block-q}
\end{equation}
and
\begin{equation}
    \bm M_{K,\Lambda}
    =
    \bm I_T
    -
    \bm Q_{K,\Lambda}
    (\bm Q_{K,\Lambda}'\bm Q_{K,\Lambda})^\dagger
    \bm Q_{K,\Lambda}'.
    \label{eq:block-m}
\end{equation}
Here \({}^\dagger\) denotes the Moore--Penrose inverse. Thus the definition
is valid for every candidate partition; when the block Gram matrix is
nonsingular, it coincides with the ordinary inverse.
The block nuisance space is
\[
    \mathcal N_{K,\Lambda}
    =
    \left\{
      \sum_{j=0}^q
      \ind(t\in I_j(\Lambda))\bm P_K(W_{t-1})'\bm c_j:
      \bm c_j\in\R^K
    \right\}.
\]
The matrix \(\bm M_{K,\Lambda}\) is the orthogonal residual-maker for this empirical space.

For each \(\beta\) and \(\Lambda\), define the profiled residual, regime-orthogonalized weight, and scalar score by
\begin{align}
    \widehat{\bm u}(\beta,\Lambda)
    &=
    \bm M_{K,\Lambda}(\bm Y-\beta\bm X_{lag}),
    \label{eq:block-profile-resid}\\
    \bm w^c(\Lambda)
    &=
    \bm M_{K,\Lambda}\bm w,
    \label{eq:block-weight}\\
    \widehat Z_t(\beta,\Lambda)
    &=
    \widehat u_t(\beta,\Lambda)w_t^c(\Lambda).
    \label{eq:block-score}
\end{align}
The empirical likelihood ratio is
\begin{align}
    L_T(\beta,\Lambda)
    &=
    \sup
    \left\{
    \prod_{t=1}^T(T\pi_t):
    \pi_t\geq0,\,
    \sum_{t=1}^T\pi_t=1,\,
    \sum_{t=1}^T\pi_t\widehat Z_t(\beta,\Lambda)=0
    \right\},
    \label{eq:block-el}\\
    \ell_T(\beta,\Lambda)
    &=
    -2\log L_T(\beta,\Lambda)
    =
    2\sum_{t=1}^T
    \log\{1+\widehat\lambda(\beta,\Lambda)\widehat Z_t(\beta,\Lambda)\},
    \label{eq:block-elr}
\end{align}
where \(\widehat\lambda(\beta,\Lambda)\) solves the scalar multiplier equation.

All minimizers below are understood as measurable selections with a fixed deterministic tie-breaking rule.

For an interval \([s,e]\), define the null-imposed segment cost
\begin{equation}
    \mathcal C_K(s,e;\beta)
    =
    \min_{\bm c\in\R^K}
    \sum_{t=s}^e
    \{Y_t-\beta X_{t-1}-\bm P_K(W_{t-1})'\bm c\}^2.
    \label{eq:segment-cost}
\end{equation}
For \(\Lambda=(k_1,\ldots,k_q)\), define
\begin{equation}
    RSS_K(\beta,\Lambda)
    =
    \sum_{j=0}^{q}
    \mathcal C_K(k_j+1,k_{j+1};\beta).
    \label{eq:rss}
\end{equation}
Let \(\mathcal L_q(\epsilon)\) be the set of \(q\)-break partitions satisfying minimum spacing. The profile partition estimator is
\begin{equation}
    \widehat\Lambda_q(\beta)
    =
    \arg\min_{\Lambda\in\mathcal L_q(\epsilon)}
    RSS_K(\beta,\Lambda).
    \label{eq:lambda-hat-q}
\end{equation}
If the number of breaks is unknown, define
\begin{equation}
    \widehat q(\beta)
    =
    \arg\min_{0\leq q\leq\bar q}
    \left\{
    RSS_K(\beta,\widehat\Lambda_q(\beta))
    +
    \operatorname{pen}_T(q,K)
    \right\},
    \label{eq:qhat}
\end{equation}
and set
\[
    \widehat\Lambda(\beta)=
    \widehat\Lambda_{\widehat q(\beta)}(\beta).
\]
The break-adaptive statistic is
\begin{equation}
    \ell_T^{BA}(\beta)
    =
    \ell_T\{\beta,\widehat\Lambda(\beta)\}.
    \label{eq:ba-stat}
\end{equation}

\section{Assumptions}

\begin{assumption}[Finite nuisance breaks and smooth regimes]
\label{ass:breaks}
The number \(q_0\) is fixed. For some \(\epsilon>0\),
\[
    \min_{0\leq j\leq q_0}(k_{j+1,0}-k_{j,0})\geq \epsilon T.
\]
Each \(m_j\) belongs to the same smoothness class \(\mathcal H^s\). If \(q_0\geq1\), the adjacent breaks satisfy
\[
    \Delta_T
    =
    \min_{0\leq j<q_0}
    \|m_{j+1}-m_j\|_{L^2(P_W)}
    >0,
\]
with \(T\Delta_T^2\) large enough relative to the block sieve dimension and logarithmic search complexity. A sufficient fixed-break version is \(T\Delta_T^2/(K\log T)\to\infty\). When \(q_0=0\), this condition is void.
\end{assumption}

\begin{assumption}[Residualized-weight nondegeneracy]
\label{ass:weight-nondegenerate}
For the true partition,
\[
    T^{-1}\bm w'\bm M_{K,\Lambda_0}\bm w
    \xrightarrow{p} Q_w,
    \qquad
    P(Q_w>\underline q)=1
\]
for some \(\underline q>0\).
\end{assumption}

\begin{assumption}[Block-sieve rates and projection bias]
\label{ass:block-sieve}
For the true partition and each regime \(j\), there exists \(\bm c_{j,K}\) such that
\[
    m_j(W_{t-1})
    =
    \bm P_K(W_{t-1})'\bm c_{j,K}
    +
    r_{j,K,t},
\]
with \(\max_j\|\bm r_{j,K}\|_T=O_p(a_K)\) and \(\max_j\|\bm r_{j,K}\|_\infty=O_p(a_K)\). Let \(\bm r_0\) denote the stacked approximation-error vector under \(\Lambda_0\). The directional projection bias satisfies
\[
    T^{-1/2}\bm r_0'\bm M_{K,\Lambda_0}\bm w=o_p(1).
\]
The block Gram matrix \(T^{-1}\bm Q_{K,\Lambda_0}'\bm Q_{K,\Lambda_0}\) has eigenvalues bounded away from zero and infinity with probability tending to one. If \(\zeta_K=\max_t\|\bm P_K(W_{t-1})\|\), then
\[
    a_K(1+\zeta_K)\to0,
    \qquad
    a_K(1+\zeta_K)^2=o(\sqrt T),
\]
and
\[
    \frac{(q_0+1)K}{\sqrt T}\to0,
    \qquad
    \zeta_K\sqrt{\frac{(q_0+1)K}{T}}\to0.
\]
\end{assumption}

\begin{remark}[Projection bias versus undersmoothing]
The condition \(T^{-1/2}\bm r_0'\bm M_{K,\Lambda_0}\bm w=o_p(1)\) only requires the nuisance approximation error to be negligible in the residualized score direction. It is weaker than the full root-\(T\) undersmoothing condition \(\sqrt T a_K\to0\). A sufficient primitive route is a product-rate condition \(\sqrt T\,a_{m,K}a_{w,K}\to0\), where \(a_{m,K}\) approximates the nuisance functions and \(a_{w,K}\) approximates the conditional mean of the weight \(w_t\) given \(W_{t-1}\).
\end{remark}

\begin{assumption}[Known-partition oracle score]
\label{ass:oracle}
With \(\bm v_0=\bm M_{K,\Lambda_0}\bm w\), define \(\widetilde{\bm U}_0=\bm M_{K,\Lambda_0}\bm U\) and
\[
    A_T^{(q)}
    =
    T^{-1/2}\sum_{t=1}^T U_tv_{t,0},
    \qquad
    V_T^{(q)}
    =
    T^{-1}\sum_{t=1}^T U_t^2v_{t,0}^2.
\]
There exist a random scale \(\eta_q\), a constant \(c_\eta>0\), and a standard normal \(Z\), independent of the limiting sigma-field generating \(\eta_q\), such that
\[
    A_T^{(q)}\xrightarrow{\mathrm{st}}\eta_qZ,
    \qquad
    V_T^{(q)}\xrightarrow{p}\eta_q^2,
    \qquad
    P(\eta_q\geq\sqrt{c_\eta})=1.
\]
In addition,
\[
    \max_{t\leq T}|U_tv_{t,0}|=o_p(\sqrt T),
\]
\[
    T^{-1}\sum_{t=1}^T
    \{(\widetilde U_{0,t}v_{t,0})^2-(U_tv_{t,0})^2\}
    =o_p(1),
    \qquad
    \max_{t\leq T}|(\widetilde U_{0,t}-U_t)v_{t,0}|
    =o_p(\sqrt T).
\]
Finally, the feasible true-partition score array satisfies the scalar empirical-likelihood convex-hull condition
\[
    P\left\{
    \min_{t\leq T}\widehat Z_t(\beta_0,\Lambda_0)<0<
    \max_{t\leq T}\widehat Z_t(\beta_0,\Lambda_0)
    \right\}\to1.
\]
This condition is imposed on the feasible array \(\widehat Z_t(\beta_0,\Lambda_0)\), not only on the oracle array \(U_tv_{t,0}\).
\end{assumption}

For the fixed-order analysis, define
\[
    \widetilde\Lambda
    =
    \widehat\Lambda_{q_0}(\beta_0),
    \qquad
    \widetilde\Lambda=(\widetilde k_1,\ldots,\widetilde k_{q_0}).
\]
The unknown-order estimator \(\widehat\Lambda(\beta_0)\) is used only in Corollary \ref{cor:qhat}.

\begin{assumption}[Fixed-order partition localization]
\label{ass:partition-localization}
For the true number of breaks,
\[
    \max_{1\leq j\leq q_0}
    |\widetilde k_j-k_{j,0}|
    =
    O_p(r_T),
    \qquad
    r_T=o(\sqrt T).
\]
Equivalently, the number of fixed-order observations assigned to the wrong regime satisfies
\[
    |\mathcal M_T|
    =
    \left|
    \{t:j_{\widetilde\Lambda}(t)\neq j_{\Lambda_0}(t)\}
    \right|
    =
    o_p(\sqrt T).
\]
This localization condition does not by itself imply score equivalence, because moving boundary observations also perturbs the fitted regime coefficients outside \(\mathcal M_T\).
\end{assumption}

\begin{assumption}[High-level fixed-order estimated-score equivalence]
\label{ass:estimated-score}
Under \(H_0:\beta=\beta_0\),
\begin{align}
    T^{-1/2}\sum_{t=1}^T
    \{\widehat Z_t(\beta_0,\widetilde\Lambda)-\widehat Z_t(\beta_0,\Lambda_0)\}
    &=o_p(1),
    \label{eq:estimated-score-num}\\
    T^{-1}\sum_{t=1}^T
    \{\widehat Z_t(\beta_0,\widetilde\Lambda)^2-\widehat Z_t(\beta_0,\Lambda_0)^2\}
    &=o_p(1),
    \label{eq:estimated-score-den}\\
    \max_{t\leq T}|\widehat Z_t(\beta_0,\widetilde\Lambda)|&=o_p(\sqrt T).
    \label{eq:estimated-score-max}
\end{align}
Zero lies in the interior of the scalar empirical-likelihood convex hull for the fixed-order estimated-partition score array with probability tending to one. This is a high-level score-equivalence condition; Assumption \ref{ass:partition-localization} alone is not claimed to prove it.
\end{assumption}

\begin{remark}[Estimated-partition bridge as the main proof risk]
The critical gap between the oracle Wilks theorem and the feasible break-adaptive statistic is not the finite-sample projection algebra. It is the bridge
\[
    S_T(\widetilde\Lambda)-S_T(\Lambda_0)=o_p(1),
    \qquad
    V_T(\widetilde\Lambda)-V_T(\Lambda_0)=o_p(1),
\]
where \(S_T(\Lambda)\) and \(V_T(\Lambda)\) are defined below. Assumption \ref{ass:partition-localization} controls how many boundary observations are assigned to the wrong regime, but it does not by itself control the change in fitted block coefficients away from the boundary. Assumption \ref{ass:estimated-score} is therefore kept as a high-level bridge unless a separate perturbation argument establishes these two displays.
\end{remark}

\section{Statements and Proofs}

The stable no-break score is
\[
    S_T^{\mathrm{st}}(\beta_0)
    =
    T^{-1/2}(\bm Y-\beta_0\bm X_{lag})'\bm M_{K,0}\bm w,
\]
where \(0\) denotes the no-break partition.

\begin{proposition}[Stable-score decomposition]
\label{prop:stable-decomp}
Under \eqref{eq:model}, the stable-nuisance score satisfies the exact identity
\[
    S_T^{\mathrm{st}}(\beta_0)
    =
    A_T^{\mathrm{st}}
    +
    B_T,
\]
where
\[
    A_T^{\mathrm{st}}
    =
    T^{-1/2}\bm U'\bm M_{K,0}\bm w,
    \qquad
    B_T
    =
    T^{-1/2}\bm\mu'\bm M_{K,0}\bm w.
\]
The term \(B_T\) is the stable-projection nuisance component left by imposing the no-break nuisance projection. When \(q_0=0\), this component is asymptotically negligible under the projection-bias condition in Assumption \ref{ass:block-sieve}; when \(q_0>0\), a nonnegligible component is interpreted as break-induced.
\end{proposition}

\begin{proof}
Under \(H_0:\beta=\beta_0\),
\[
    \bm Y-\beta_0\bm X_{lag}=\bm\mu+\bm U.
\]
Multiplying by the stable residualized weight gives
\[
    T^{-1/2}(\bm Y-\beta_0\bm X_{lag})'\bm M_{K,0}\bm w
    =
    T^{-1/2}\bm U'\bm M_{K,0}\bm w
    +
    T^{-1/2}\bm\mu'\bm M_{K,0}\bm w.
\]
This is an exact algebraic identity. No stable-sieve approximation remainder is needed because \(B_T\) is defined using the exact nuisance vector \(\bm\mu\).

When \(q_0=0\), write \(\bm\mu=\bm P\bm c_K+\bm r_K\). Since \(\bm M_{K,0}\bm P\bm c_K=\bm0\),
\[
    B_T
    =
    T^{-1/2}\bm r_K'\bm M_{K,0}\bm w
    =o_p(1)
\]
by the projection-bias condition in Assumption \ref{ass:block-sieve} with \(\Lambda_0=0\).
\end{proof}

\begin{definition}[Break-induced spurious predictivity]
The stable nuisance specification generates break-induced spurious predictivity at \(\beta_0\) if
\[
    T^{-1/2}
    \langle \bm M_{K,0}\bm\mu,\bm M_{K,0}\bm w\rangle_{\R^T}
    =
    B_T
    \not=o_p(1).
\]
The condition is directional: a nuisance break with little projection on the residualized persistent weight need not distort the test.
\end{definition}

For the stable no-break empirical-likelihood statistic, write
\[
    V_T^{\mathrm{st}}
    =
    T^{-1}\sum_{t=1}^T \widehat Z_t(\beta_0,0)^2,
\]
where \(\widehat Z_t(\beta_0,0)\) is the score in \eqref{eq:block-score} computed under the no-break partition.

\begin{theorem}[Omitted-break failure]
\label{thm:failure}
Suppose the stable empirical-likelihood statistic admits
\[
    \ell_T^{\mathrm{st}}(\beta_0)
    =
    \frac{\{S_T^{\mathrm{st}}(\beta_0)\}^2}{V_T^{\mathrm{st}}}
    +o_p(1),
\]
with \(V_T^{\mathrm{st}}\) bounded away from zero in probability. Define
\[
    a_T=\frac{A_T^{\mathrm{st}}}{(V_T^{\mathrm{st}})^{1/2}},
    \qquad
    b_T=\frac{B_T}{(V_T^{\mathrm{st}})^{1/2}}.
\]
If \(a_T=O_p(1)\) and \(|b_T|\to_p\infty\), then, for every fixed chi-square critical value \(c_\alpha\),
\[
    P\{\ell_T^{\mathrm{st}}(\beta_0)>c_\alpha\}\to1.
\]
If instead
\[
    (a_T,b_T)\xrightarrow{\mathrm{st}}(Z,B),
\]
where \(Z\sim N(0,1)\) is independent of the limiting sigma-field and \(B\) is measurable with respect to that sigma-field, then
\[
    \ell_T^{\mathrm{st}}(\beta_0)\Rightarrow (Z+B)^2.
\]
Conditionally on \(B\), the limit is \(\chi_1^2(B^2)\). Thus a constant \(B=\delta\) gives \(\chi_1^2(\delta^2)\), a random \(B\) gives a mixture of noncentral chi-squares, and \(B=0\) almost surely gives the central \(\chi_1^2\) limit.
\end{theorem}

\begin{proof}
By Proposition \ref{prop:stable-decomp},
\[
    S_T^{\mathrm{st}}(\beta_0)=A_T^{\mathrm{st}}+B_T.
\]
Hence the scalar expansion is \(\ell_T^{\mathrm{st}}(\beta_0)=(a_T+b_T)^2+o_p(1)\). If \(a_T=O_p(1)\) and \(|b_T|\to_p\infty\), then
\[
    |a_T+b_T|\geq |b_T|-|a_T|\to_p\infty,
\]
so the statistic diverges in probability. Under joint stable convergence, the continuous mapping theorem gives \((a_T+b_T)^2\Rightarrow (Z+B)^2\). Since \(Z\) is conditionally standard normal given the limiting sigma-field, \((Z+B)^2\mid B\sim\chi_1^2(B^2)\).
\end{proof}

\begin{remark}[Scope of omitted-break divergence]
The divergence part of Theorem \ref{thm:failure} is conditional on the displayed stable empirical-likelihood quadratic expansion. With a fixed omitted nuisance break, \(B_T\) can be of order \(\sqrt T\), and the null-style scalar expansion need not follow automatically. Establishing generic fixed-break divergence therefore requires a separate expansion under the misspecified stable nuisance projection.
\end{remark}

\begin{lemma}[Block projection identities]
\label{lem:block-identities}
For any partition \(\Lambda\),
\[
    \bm Q_{K,\Lambda}'\bm M_{K,\Lambda}=\bm0,
    \qquad
    \bm M_{K,\Lambda}'=\bm M_{K,\Lambda},
    \qquad
    \bm M_{K,\Lambda}^2=\bm M_{K,\Lambda}.
\]
For every \(\bm c\) conformable with \(\bm Q_{K,\Lambda}\),
\[
    (\bm Q_{K,\Lambda}\bm c)'\bm M_{K,\Lambda}\bm w=0.
\]
\end{lemma}

\begin{proof}
Let
\[
    \bm\Pi_{K,\Lambda}
    =
    \bm Q_{K,\Lambda}(\bm Q_{K,\Lambda}'\bm Q_{K,\Lambda})^\dagger\bm Q_{K,\Lambda}'.
\]
This is the orthogonal projector onto the column space of \(\bm Q_{K,\Lambda}\). Hence \(\bm\Pi_{K,\Lambda}'=\bm\Pi_{K,\Lambda}\), \(\bm\Pi_{K,\Lambda}^2=\bm\Pi_{K,\Lambda}\), and \(\bm Q_{K,\Lambda}'(\bm I_T-\bm\Pi_{K,\Lambda})=\bm0\). Since \(\bm M_{K,\Lambda}=\bm I_T-\bm\Pi_{K,\Lambda}\), the displayed identities follow. The final equality is obtained by multiplying \(\bm Q_{K,\Lambda}'\bm M_{K,\Lambda}=\bm0\) by \(\bm c\) and \(\bm w\).
\end{proof}

Set \(\Lambda=\Lambda_0\) and write
\[
    \bm v_0=\bm M_{K,\Lambda_0}\bm w,
    \qquad
    \widehat{\bm u}_0(\beta_0)=\bm M_{K,\Lambda_0}(\bm Y-\beta_0\bm X_{lag}).
\]
Under the true partition,
\[
    \bm Y-\beta_0\bm X_{lag}
    =
    \bm Q_{K,\Lambda_0}\bm c_0+\bm r_0+\bm U,
\]
where \(\bm r_0\) stacks the regime-specific sieve approximation errors. Let \(\bm e_0=\bm M_{K,\Lambda_0}\bm r_0\).

\begin{lemma}[Known-partition replacement]
\label{lem:known-replacement}
Under Assumptions \ref{ass:block-sieve} and \ref{ass:oracle},
\begin{align}
    T^{-1/2}\sum_{t=1}^T\widehat Z_t(\beta_0,\Lambda_0)
    &=T^{-1/2}\sum_{t=1}^T U_tv_{t,0}+o_p(1),
    \label{eq:known-num}\\
    T^{-1}\sum_{t=1}^T\widehat Z_t(\beta_0,\Lambda_0)^2
    &=T^{-1}\sum_{t=1}^T U_t^2v_{t,0}^2+o_p(1),
    \label{eq:known-den}\\
    \max_{t\leq T}|\widehat Z_t(\beta_0,\Lambda_0)|&=o_p(\sqrt T).
    \label{eq:known-max}
\end{align}
\end{lemma}

\begin{proof}
Using Lemma \ref{lem:block-identities},
\[
    \widehat{\bm u}_0(\beta_0)
    =
    \bm M_{K,\Lambda_0}(\bm U+\bm r_0),
    \qquad
    \bm v_0=\bm M_{K,\Lambda_0}\bm w.
\]
Because \(\bm M_{K,\Lambda_0}\) is symmetric and idempotent,
\[
    T^{-1/2}\widehat{\bm u}_0(\beta_0)'\bm v_0
    =
    T^{-1/2}(\bm U+\bm r_0)'\bm v_0.
\]
By Assumption \ref{ass:block-sieve},
\[
    T^{-1/2}\bm r_0'\bm v_0
    =
    T^{-1/2}\bm r_0'\bm M_{K,\Lambda_0}\bm w
    =o_p(1),
\]
which proves \eqref{eq:known-num}.

It remains to record the denominator and maximum bounds. Let \(G_T=T^{-1}\bm Q_{K,\Lambda_0}'\bm Q_{K,\Lambda_0}\). Since the rows of \(\bm Q_{K,\Lambda_0}\) have norm at most \(\zeta_K\) and the eigenvalues of \(G_T\) are bounded away from zero and infinity,
\[
    \|\bm v_0\|_\infty=O_p(1+\zeta_K).
\]
Moreover,
\[
    T^{-1}\bm Q_{K,\Lambda_0}'\bm r_0=O_p(a_K),
\]
and hence
\[
    \|\bm e_0\|_T\leq\|\bm r_0\|_T=O_p(a_K),
    \qquad
    \|\bm e_0\|_\infty=O_p\{a_K(1+\zeta_K)\}.
\]
The stated rates imply
\[
    \|\bm e_0\odot\bm v_0\|_T
    =O_p\{a_K(1+\zeta_K)\}=o_p(1),
    \qquad
    \|\bm e_0\odot\bm v_0\|_\infty
    =O_p\{a_K(1+\zeta_K)^2\}=o_p(\sqrt T).
\]
For the denominator, write \(\widetilde{\bm U}_0=\bm M_{K,\Lambda_0}\bm U\). Then
\[
    \widehat Z_t(\beta_0,\Lambda_0)
    =
    \widetilde U_{0,t}v_{t,0}+e_{0,t}v_{t,0}.
\]
Assumption \ref{ass:oracle} gives
\[
    T^{-1}\sum_{t=1}^T
    \{(\widetilde U_{0,t}v_{t,0})^2-(U_tv_{t,0})^2\}=o_p(1),
\]
and Cauchy--Schwarz with \(\|\bm e_0\odot\bm v_0\|_T=o_p(1)\) gives the remaining cross and square terms. This proves \eqref{eq:known-den}. The maximum bound follows from
\[
    |\widehat Z_t(\beta_0,\Lambda_0)|
    \leq
    |U_tv_{t,0}|
    +
    |(\widetilde U_{0,t}-U_t)v_{t,0}|
    +
    |e_{0,t}v_{t,0}|,
\]
Assumption \ref{ass:oracle}, and \(\|\bm e_0\odot\bm v_0\|_\infty=o_p(\sqrt T)\).
\end{proof}

\begin{lemma}[Scalar empirical-likelihood expansion]
\label{lem:el-expansion}
Let \(Z_{t,T}\) be scalar scores and \(\widehat V_T^Z=T^{-1}\sum_tZ_{t,T}^2\). Suppose zero lies in the interior of the scalar convex hull with probability tending to one and
\[
    T^{-1/2}\sum_t Z_{t,T}=O_p(1),
    \qquad
    \widehat V_T^Z\to_p\eta^2,
    \qquad
    \max_t|Z_{t,T}|=o_p(\sqrt T),
\]
where \(P(\eta^2\geq c)=1\) for some fixed \(c>0\). Then
\[
    2\sum_{t=1}^T\log(1+\widehat\lambda Z_{t,T})
    =
    \frac{\{T^{-1/2}\sum_tZ_{t,T}\}^2}
         {T^{-1}\sum_tZ_{t,T}^2}
    +o_p(1).
\]
\end{lemma}

\begin{proof}
The second-moment condition implies \(P(\widehat V_T^Z\geq c/2)\to1\). Let \(S_T^Z=\sum_tZ_{t,T}\) and \(Q_T^Z=\sum_tZ_{t,T}^2\). On this event, \(S_T^Z=O_p(\sqrt T)\) and \(Q_T^Z\geq cT/2\). Monotonicity of the scalar multiplier equation, evaluated at \(\pm 2|S_T^Z|/Q_T^Z\), gives
\[
    |\widehat\lambda|
    \leq
    2|S_T^Z|/Q_T^Z
    =O_p(T^{-1/2}).
\]
Since \(\max_t|\widehat\lambda Z_{t,T}|=o_p(1)\), a Taylor expansion of the multiplier equation gives
\[
    \widehat\lambda
    =
    \frac{\sum_t Z_{t,T}}{\sum_t Z_{t,T}^2}
    +o_p(T^{-1/2}).
\]
Substitution into the second-order expansion of the log empirical likelihood gives the displayed quadratic approximation.
\end{proof}

\begin{theorem}[Known-partition Wilks theorem]
\label{thm:known-wilks}
Under Assumptions \ref{ass:breaks}, \ref{ass:weight-nondegenerate}, \ref{ass:block-sieve}, and \ref{ass:oracle}, under \(H_0:\beta=\beta_0\),
\[
    \ell_T(\beta_0,\Lambda_0)\Rightarrow\chi_1^2.
\]
\end{theorem}

\begin{proof}
Lemma \ref{lem:known-replacement} reduces the feasible score to the matched block-projection oracle score \(U_tv_{t,0}\). Assumption \ref{ass:oracle} gives
\[
    A_T^{(q)}\xrightarrow{\mathrm{st}}\eta_qZ,
    \qquad
    V_T^{(q)}\to_p\eta_q^2,
    \qquad
    P(\eta_q\geq\sqrt{c_\eta})=1.
\]
Stable Slutsky therefore gives \(A_T^{(q)}/\sqrt{V_T^{(q)}}\Rightarrow Z\). Lemma \ref{lem:el-expansion} yields
\[
    \ell_T(\beta_0,\Lambda_0)
    =
    \frac{\{T^{-1/2}\sum_tU_tv_{t,0}\}^2}
         {T^{-1}\sum_tU_t^2v_{t,0}^2}
    +o_p(1)
    \Rightarrow Z^2\sim\chi_1^2.
\]
\end{proof}

\[
    S_T(\Lambda)
    =
    T^{-1/2}\sum_{t=1}^T\widehat Z_t(\beta_0,\Lambda),
    \qquad
    V_T(\Lambda)
    =
    T^{-1}\sum_{t=1}^T\widehat Z_t(\beta_0,\Lambda)^2.
\]
By symmetry and idempotence of \(\bm M_{K,\Lambda}\),
\begin{equation}
    S_T(\Lambda)
    =
    T^{-1/2}
    (\bm Y-\beta_0\bm X_{lag})'\bm M_{K,\Lambda}\bm w.
    \label{eq:score-numerator-projection}
\end{equation}

\begin{lemma}[Partition perturbation implies score perturbation: crude route]
\label{lem:partition-perturbation}
Let
\[
    \Delta_M
    =
    \bm M_{K,\widetilde\Lambda}-\bm M_{K,\Lambda_0}.
\]
Consider the fixed-order case with the correct number of breaks and
\[
    \max_j|\widetilde k_j-k_{j,0}|=O_p(r_T),
    \qquad
    r_T=o(\sqrt T).
\]
Suppose bounded weights, bounded nuisance functions, bounded basis envelope \(\zeta_K\), strong minimum spacing, and uniform Gram stability over the \(r_T\)-neighborhood of \(\Lambda_0\) imply the perturbation bounds
\begin{align}
    T^{-1/2}\bm U'\Delta_M\bm w
    &=
    O_p\!\left(\zeta_K\sqrt{\frac{K r_T}{T}}\right)+o_p(1),
    \label{eq:bridge-u-bound}\\
    T^{-1/2}\bm\mu'\Delta_M\bm w
    &=
    O_p\!\left(\frac{(1+\zeta_K)r_T}{\sqrt T}\right)+o_p(1),
    \label{eq:bridge-mu-bound}\\
    V_T(\widetilde\Lambda)-V_T(\Lambda_0)
    &=
    O_p\!\left(
      \frac{(1+\zeta_K)^2r_T}{T}
      +
      \frac{\zeta_K^2K r_T}{T}
    \right)+o_p(1),
    \label{eq:bridge-v-bound}\\
    \max_{t\leq T}|\widehat Z_t(\beta_0,\widetilde\Lambda)|
    &=o_p(\sqrt T).
    \label{eq:bridge-max-bound}
\end{align}
If
\[
    \zeta_K\sqrt{\frac{K r_T}{T}}\to0,
    \qquad
    \frac{(1+\zeta_K)r_T}{\sqrt T}\to0,
    \qquad
    \frac{\zeta_K^2K r_T}{T}\to0,
\]
then the numerator, denominator, and maximum-score parts of Assumption \ref{ass:estimated-score} hold. If the estimated-partition convex-hull condition also holds, then the full fixed-order estimated-score equivalence assumption holds.
\end{lemma}

\begin{proof}
Under \(H_0:\beta=\beta_0\), \(\bm Y-\beta_0\bm X_{lag}=\bm\mu+\bm U\). Using \eqref{eq:score-numerator-projection},
\[
    S_T(\widetilde\Lambda)-S_T(\Lambda_0)
    =
    T^{-1/2}(\bm U+\bm\mu)'\Delta_M\bm w.
\]
The two numerator bounds \eqref{eq:bridge-u-bound} and \eqref{eq:bridge-mu-bound}, together with the displayed rates, give
\[
    S_T(\widetilde\Lambda)-S_T(\Lambda_0)=o_p(1).
\]
The denominator requirement is \eqref{eq:bridge-v-bound}. The maximum-score requirement is \eqref{eq:bridge-max-bound}. The estimated-partition convex-hull condition must still be verified separately or imposed as part of Assumption \ref{ass:estimated-score}.
\end{proof}

\begin{remark}[How to use Lemma \ref{lem:partition-perturbation}]
This is intentionally a crude de-risking lemma. Its role is to isolate the proof target for a primitive one-break, known-\(q_0\), strong-spacing analysis. Boundary observations should generate the \(r_T/\sqrt T\) nuisance term; least-squares coefficient perturbation should generate the \(\zeta_K\sqrt{K r_T/T}\) stochastic term; and the denominator should be smaller after \(T^{-1}\) normalization. Proving those perturbation bounds in the simplest bounded-basis case would turn the estimated-partition bridge from a purely high-level assumption into a primitive sufficient condition.
\end{remark}

\begin{lemma}[Fixed-order estimated-partition likelihood transfer]
\label{lem:estimated-transfer}
Under the assumptions of Theorem \ref{thm:known-wilks} and Assumption \ref{ass:estimated-score},
\[
    \ell_T(\beta_0,\widetilde\Lambda)
    -
    \ell_T(\beta_0,\Lambda_0)
    =o_p(1).
\]
\end{lemma}

\begin{proof}
Lemma \ref{lem:known-replacement} and Assumption \ref{ass:oracle} give \(S_T(\Lambda_0)=O_p(1)\), \(V_T(\Lambda_0)\to_p\eta_q^2\), and a true-partition maximum-score bound. Assumption \ref{ass:estimated-score} gives
\[
    S_T(\widetilde\Lambda)-S_T(\Lambda_0)=o_p(1),
    \qquad
    V_T(\widetilde\Lambda)-V_T(\Lambda_0)=o_p(1),
\]
and the fixed-order estimated-partition maximum-score and convex-hull conditions. Lemma \ref{lem:el-expansion} applies to both arrays, and the two quadratic ratios differ by \(o_p(1)\) because both denominators are bounded away from zero with probability tending to one.
\end{proof}

\begin{theorem}[Fixed-order estimated-partition Wilks theorem]
\label{thm:estimated-equiv}
Under the assumptions of Theorem \ref{thm:known-wilks} and Assumption \ref{ass:estimated-score},
\[
    \ell_T(\beta_0,\widetilde\Lambda)\Rightarrow\chi_1^2.
\]
This theorem is high-level: it proves that fixed-order estimated-score equivalence implies empirical-likelihood equivalence, not that the profile-sieve partition estimator satisfies Assumption \ref{ass:estimated-score} from primitive conditions.
\end{theorem}

\begin{proof}
By Lemma \ref{lem:estimated-transfer},
\[
    \ell_T(\beta_0,\widetilde\Lambda)
    =
    \ell_T(\beta_0,\Lambda_0)+o_p(1).
\]
The result follows from Theorem \ref{thm:known-wilks}.
\end{proof}

\begin{corollary}[Unknown number of breaks]
\label{cor:qhat}
If \(P\{\widehat q(\beta_0)=q_0\}\to1\) and Assumption \ref{ass:estimated-score} holds on the fixed-order event, then
\[
    \ell_T^{BA}(\beta_0)\Rightarrow\chi_1^2.
\]
\end{corollary}

\begin{proof}
On the event \(\{\widehat q(\beta_0)=q_0\}\), \(\widehat\Lambda(\beta_0)=\widetilde\Lambda\). Since this event has probability tending to one, Theorem \ref{thm:estimated-equiv} gives the same Wilks limit for \(\ell_T^{BA}(\beta_0)\).
\end{proof}

\begin{theorem}[Local predictive alternatives]
\label{thm:local-power}
Let \(\beta_T=\beta_0+d_T\), where \(d_T\to0\). Maintain Assumptions \ref{ass:breaks}, \ref{ass:weight-nondegenerate}, and \ref{ass:block-sieve}. In addition, assume that the oracle-score part of Assumption \ref{ass:oracle} continues to hold under the local law \(P_{\beta_T}\):
\[
    T^{-1/2}\sum_{t=1}^TU_tv_{t,0}
    \xrightarrow{\mathrm{st}}\eta_qZ,
    \qquad
    T^{-1}\sum_{t=1}^TU_t^2v_{t,0}^2
    \xrightarrow{p}\eta_q^2,
    \qquad
    \max_{t\leq T}|U_tv_{t,0}|=o_p(\sqrt T),
\]
with \(P(\eta_q\geq\sqrt{c_\eta})=1\). Define
\[
    g_{t,T}
    =
    d_T(\bm M_{K,\Lambda_0}\bm X_{lag})_tv_{t,0}.
\]
Suppose
\begin{align}
    T^{-1/2}\sum_{t=1}^Tg_{t,T}&\xrightarrow{p}\eta_q\delta_h,
    \label{eq:local-drift}\\
    T^{-1}\sum_{t=1}^Tg_{t,T}^2&=o_p(1),
    \label{eq:local-l2}\\
    \max_{t\leq T}|g_{t,T}|&=o_p(\sqrt T),
    \label{eq:local-max}
\end{align}
with \(\delta_h\) fixed, and set \(Z_{t,T}^{\mathrm{loc}}=U_tv_{t,0}+g_{t,T}\). Under the local data-generating law \(\beta_T\), assume the feasible true-partition score satisfies
\begin{align}
    T^{-1/2}\sum_{t=1}^T
    \{\widehat Z_t(\beta_0,\Lambda_0)-Z_{t,T}^{\mathrm{loc}}\}
    &=o_p(1),
    \label{eq:local-repl-num}\\
    T^{-1}\sum_{t=1}^T
    \{\widehat Z_t(\beta_0,\Lambda_0)^2-(Z_{t,T}^{\mathrm{loc}})^2\}
    &=o_p(1),
    \label{eq:local-repl-den}\\
    \max_{t\leq T}
    |\widehat Z_t(\beta_0,\Lambda_0)-Z_{t,T}^{\mathrm{loc}}|
    &=o_p(\sqrt T).
    \label{eq:local-repl-max}
\end{align}
Also assume the feasible true-partition convex-hull condition holds under \(\beta_T\):
\[
    P_{\beta_T}\left\{
    \min_{t\leq T}\widehat Z_t(\beta_0,\Lambda_0)<0<
    \max_{t\leq T}\widehat Z_t(\beta_0,\Lambda_0)
    \right\}\to1.
\]
Assume, in addition, that the fixed-order estimated-score equivalence in Assumption \ref{ass:estimated-score} holds under \(\beta_T\), with all \(o_p\) statements evaluated under the local law. Then, when the statistic is evaluated at \(\beta_0\),
\[
    \ell_T(\beta_0,\widetilde\Lambda)
    \Rightarrow
    \chi_1^2(\delta_h^2).
\]
If \(P\{\widehat q(\beta_0)=q_0\}\to1\), the same limit holds for \(\ell_T^{BA}(\beta_0)\).
\end{theorem}

\begin{proof}
Under the local alternative,
\[
    \bm Y-\beta_0\bm X_{lag}
    =
    \bm\mu+\bm U+d_T\bm X_{lag}.
\]
At the true partition, the block projection removes the regime-specific sieve approximation to \(\bm\mu\). The numerator replacement condition \eqref{eq:local-repl-num} gives
\[
    T^{-1/2}\sum_{t=1}^T \widehat Z_t(\beta_0,\Lambda_0)
    =
    T^{-1/2}\sum_{t=1}^T U_tv_{t,0}
    +
    T^{-1/2}\sum_{t=1}^Tg_{t,T}
    +o_p(1).
\]
By \eqref{eq:local-drift}, the second term converges to \(\eta_q\delta_h\). The denominator replacement condition \eqref{eq:local-repl-den} reduces the feasible denominator to the shifted oracle denominator. Moreover,
\[
    T^{-1}\sum_t(U_tv_{t,0}+g_{t,T})^2
    =
    V_T^{(q)}
    +
    2T^{-1}\sum_tU_tv_{t,0}g_{t,T}
    +
    T^{-1}\sum_tg_{t,T}^2.
\]
The last term is \(o_p(1)\) by \eqref{eq:local-l2}; the cross term is \(o_p(1)\) by Cauchy--Schwarz and \(V_T^{(q)}=O_p(1)\). Thus the feasible denominator remains \(\eta_q^2+o_p(1)\). The maximum condition follows from the oracle maximum condition, \eqref{eq:local-max}, and \eqref{eq:local-repl-max}. The feasible convex-hull condition and Lemma \ref{lem:el-expansion} give the known-partition limit \((Z+\delta_h)^2\). Lemma \ref{lem:estimated-transfer}, applied under the local-alternative estimated-score and convex-hull conditions, transfers the result to \(\widetilde\Lambda\). If \(\widehat q(\beta_0)=q_0\) with probability tending to one, then \(\widehat\Lambda(\beta_0)=\widetilde\Lambda\) on that event, giving the unknown-order statement.
\end{proof}

\section{Proof Dependency Map and Risk Register}

This file is a math-spine document. It keeps only the model, assumptions, definitions, lemmas, theorems, proofs, and dependency checks needed to evaluate whether
\[
    \text{assumptions}
    \quad\Longrightarrow\quad
    \text{Wilks and local-power conclusions}.
\]
The intended audit order is
\[
\begin{gathered}
\boxed{\text{Algebra}}
\rightarrow
\boxed{\text{Stable-score identity}}\\
\rightarrow
\boxed{\text{Known-partition Wilks}}
\rightarrow
\boxed{\text{Estimated-partition bridge}}\\
\rightarrow
\boxed{\text{Unknown-order selection}}
\rightarrow
\boxed{\text{Failure and local alternatives}}.
\end{gathered}
\]

\subsection*{Dependency map}

\begin{center}
\begin{tabular}{>{\raggedright\arraybackslash}p{0.25\textwidth}
                >{\raggedright\arraybackslash}p{0.34\textwidth}
                >{\raggedright\arraybackslash}p{0.31\textwidth}}
\hline
Result & Direct argument in this file & High-level or remaining input \\
\hline
Projection algebra & Lemma \ref{lem:block-identities}; Moore--Penrose projection gives \(\bm M'=\bm M\), \(\bm M^2=\bm M\), and \(\bm Q'\bm M=\bm0\). & None beyond conformability of \(\bm Q_{K,\Lambda}\). \\
Stable-score decomposition & Proposition \ref{prop:stable-decomp}; exact identity \(S_T^{\mathrm{st}}=T^{-1/2}\bm U'\bm M_{K,0}\bm w+T^{-1/2}\bm\mu'\bm M_{K,0}\bm w\). & Do not overclaim fixed-break divergence without the stable EL expansion in Theorem \ref{thm:failure}. \\
Known-partition Wilks & Lemmas \ref{lem:known-replacement} and \ref{lem:el-expansion}, then Theorem \ref{thm:known-wilks}. & Oracle mixed-normal score, feasible convex hull, max condition, and projection-noise replacement in Assumption \ref{ass:oracle}. \\
Estimated fixed-order partition & Lemma \ref{lem:estimated-transfer} transfers EL once numerator, denominator, max, and convex-hull equivalence hold. & Assumption \ref{ass:estimated-score}; Lemma \ref{lem:partition-perturbation} records a crude route to primitive sufficient conditions. \\
Unknown number of breaks & Corollary \ref{cor:qhat}; reduce to fixed-order theorem on \(\{\widehat q(\beta_0)=q_0\}\). & Order-selection consistency \(P\{\widehat q(\beta_0)=q_0\}\to1\). \\
Local alternatives & Theorem \ref{thm:local-power}; shifted score \(U_tv_{t,0}+g_{t,T}\) gives \(\chi_1^2(\delta_h^2)\). & Oracle score limit under the local law, local shifted-array replacement, convex hull, and estimated-score equivalence. \\
\hline
\end{tabular}
\end{center}

\subsection*{Risk register}

\begin{center}
\begin{tabular}{>{\raggedright\arraybackslash}p{0.26\textwidth}
                >{\raggedright\arraybackslash}p{0.24\textwidth}
                >{\raggedright\arraybackslash}p{0.11\textwidth}
                >{\raggedright\arraybackslash}p{0.29\textwidth}}
\hline
Component & Current status & Risk & Next mathematical action \\
\hline
Projection algebra & Exact finite-sample result. & Low & Keep the Moore--Penrose formulation so candidate singular Gram matrices are covered. \\
Stable-score decomposition & Exact identity. & Low & Maintain the distinction between score contamination and automatic fixed-break EL divergence. \\
Known-partition Wilks & Conditional theorem is coherent. & Medium & Verify the oracle CLT, denominator, full-sample projection, and max condition from primitive persistence assumptions. \\
Estimated-partition bridge & High-level transfer plus crude perturbation target. & High & Prove Lemma \ref{lem:partition-perturbation}'s bounds first for one break, known \(q_0\), bounded basis, strong spacing, and martingale or i.i.d. errors. \\
Unknown \(q\) & Conditional corollary. & Medium--high & Keep \(P\{\widehat q(\beta_0)=q_0\}\to1\) as a high-level condition until the fixed-order bridge is de-risked. \\
Omitted-break failure & Conditional on stable EL expansion. & Medium & Do not state generic fixed-break divergence without a misspecified stable-EL expansion. \\
Local alternatives & Conditional shifted-score theorem. & Medium & Check the local-law oracle score and replacement relative to \(U_tv_{t,0}+g_{t,T}\), not only relative to \(U_tv_{t,0}\). \\
\hline
\end{tabular}
\end{center}

\subsection*{Exit criteria for paper polish}

The theory is ready to support story-level rewriting only after the following three questions have precise answers:
\begin{enumerate}
    \item Why does the known-partition statistic have a \(\chi^2_1\) limit?
    \item Why does replacing \(\Lambda_0\) by \(\widetilde\Lambda\) not change the score to first order?
    \item Is unknown \(q_0\) proved from primitive penalty conditions, or is it imposed through \(P\{\widehat q(\beta_0)=q_0\}\to1\)?
\end{enumerate}
\end{document}